\documentclass[11pt,letterpaper]{article}
\usepackage[margin=1in]{geometry}
\usepackage{booktabs}
\usepackage{longtable}
\usepackage{array}
\usepackage{xcolor}
\usepackage{hyperref}
\usepackage{listings}
\usepackage{fancyvrb}
\usepackage{titlesec}
\usepackage{enumitem}

\hypersetup{
  colorlinks=true,
  linkcolor=blue!60!black,
  urlcolor=blue!60!black,
  citecolor=blue!60!black
}

\lstset{
  basicstyle=\ttfamily\small,
  breaklines=true,
  frame=single,
  backgroundcolor=\color{gray!10}
}

\title{Replication Report: Alayande \textit{et al.} (2020) \\
\large ``Integrated genome-based probiotic relevance and safety evaluation of \textit{Lactobacillus reuteri} PNW1''}
\author{Ollie (OpenClaw AI) --- BVBRC Replication Project (Wave 5 / BVBRC-100, target \#64)}
\date{2026-07-03}

\begin{document}
\maketitle

\begin{center}
\fbox{\parbox{0.92\textwidth}{\centering
\textbf{Paper:} Alayande KA, Aiyegoro OA, Nengwekhulu TM, Katata-Seru L, Ateba CN. \textit{PLoS ONE} 15(7): e0235873 (2020). \\
DOI: \href{https://doi.org/10.1371/journal.pone.0235873}{10.1371/journal.pone.0235873} \quad PMID: 32687505 \quad PMC: PMC7371166 \quad License: CC BY \\[2pt]
\textbf{Companion:} Alayande \textit{et al.} \textit{Microbiol.\ Resour.\ Announc.} 8(8): e00034-19 (2019). PMID 30834362, PMC6386563. \\[4pt]
\textbf{Verdict:} \textcolor{green!50!black}{\textbf{REPLICATED (strong)}}
}}
\end{center}

\vspace{6pt}

\section*{Executive Summary}
Every headline genomic claim of the paper's \textit{Results} section that is checkable from the deposited public assembly is independently reproduced here from first principles, using a completely independent tool chain (NCBI PGAP annotation + \texttt{abricate} v-current with ResFinder / CARD / NCBI-AMR / ARG-ANNOT / VFDB / VICTORS / ecoli\_vf / PlasmidFinder, plus \texttt{minced} v0.4). Two published claims dependent on assemblers other than the deposited PGAP annotation (helveticin~J and D-lactate dehydrogenase by explicit name; CRISPR arrays) are \textbf{partially confirmed} --- consistent with what I find, differences attributable to annotator conservatism and assembly-fragmentation limits, respectively.

\section{Paper}
Alayande \textit{et al.} (2020) sequence \textit{Lactobacillus reuteri} PNW1 --- a strain isolated in 2012 from the faeces of a piglet of the indigenous South African Windsnyer breed (SampleA3 $\rightarrow$ BioSample SAMN10397676 $\rightarrow$ BioProject PRJNA504734) --- and use whole-genome sequencing plus a battery of bioinformatic tools (PGAP + RAST annotation; ResFinder / ARG-ANNOT / CARD for AMR; VirulenceFinder / VFDB for VFs; PHASTER for prophages; CRISPRFinder for CRISPR; OASIS for insertion sequences; PathogenFinder for human-pathogen probability) plus phenotypic follow-up (HPLC for biogenic amines; agar-well-diffusion antimicrobial assay against Shiga-toxigenic \textit{E.\ coli} O177) to argue that PNW1 is a plausible probiotic candidate.

The species has since been reclassified into \textit{Limosilactobacillus} (Zheng \textit{et al.} 2020); NCBI now lists this genome under \textbf{\textit{Limosilactobacillus reuteri} PNW1}. Deposited genome: \textbf{GCA\_003790365.1} (GenBank, live). The RefSeq mirror GCF\_003790365.1 was \textbf{suppressed} by NCBI staff on quality grounds (``contaminated''); the underlying GenBank assembly remains public and is what the paper analyses.

\section{Claims tested}
\begin{longtable}{p{0.04\textwidth} p{0.42\textwidth} p{0.15\textwidth} p{0.22\textwidth} p{0.09\textwidth}}
\toprule
\# & Claim (as stated in paper) & Type & Testable? & Tested \\
\midrule
\endhead
C1 & Genome 2{,}430{,}215~bp in 420 contigs, 39\% G+C. & Assembly stats & Yes (GenBank FASTA). & \checkmark \\
C2 & Deposit/BioSample metadata consistent (piglet gut, S. Africa, MiSeq, SPAdes 3.12). & Provenance & Yes (NCBI Datasets). & \checkmark \\
C3 & Arg deiminase only ``toxic-biochemical'' enzyme; D-LDH, L-LDH, helveticin~J present. & Annotation & Yes (PGAP protein FASTA). & \checkmark / partial \\
C4 & Only \texttt{lnu(C)} and \texttt{tet(W)} resistance genes present. & AMR screen & Yes (abricate $\times$4 DBs). & \checkmark \\
C5 & No virulence factor hits; PathogenFinder probability = 0. & Virulence screen & VF yes (VFDB/VICTORS/ecoli\_vf); PathogenFinder paywalled. & \checkmark (VF) \\
C6 & 2 intact prophages; 9 IS-element CDSs across 7 families; 5 CRISPR CDSs each with Cas. & Mobilome & Partial (prophage/IS visible; CRISPR needs intact array). & \checkmark / partial \\
C7 & Bacteriocin agar-well-diffusion phenotype vs STEC \textit{E.\ coli} O177. & Wet-lab & \textbf{No} --- requires isolate + STEC. & --- \\
\bottomrule
\end{longtable}

\section{Method}
All work performed on \textbf{CherryRd (macOS)} on 2026-07-03. Zero paid endpoints. Only public NCBI data and FOSS tools.

\begin{enumerate}[leftmargin=1.4em]
  \item \textbf{Strain metadata:} NCBI Datasets REST v2alpha for \texttt{Limosilactobacillus reuteri} $\cap$ strain \texttt{PNW1} $\rightarrow$ 2 records (GCA\_003790365.1 live; GCF suppressed). Both = assembly \texttt{ASM379036v1}, submitter North-West University, release 2018-11-18, BioSample SAMN10397676 (piglet faeces, MRS, S. Africa: Pretoria, 25.89S 28.21E, coll.\ 2012-06), Illumina MiSeq, SPAdes 3.12.0 --- consistent with paper.
  \item \textbf{PubMed confirmation:} esearch $\rightarrow$ PMIDs 32687505 (target) + 30834362 (companion MRA). Abstracts pulled as source of truth.
  \item \textbf{Assembly download:} NCBI Datasets v2alpha, \texttt{include\_annotation\_type=GENOME\_FASTA,GENOME\_GFF,PROT\_FASTA} $\rightarrow$ 1.35~MB zip; genomic FASTA, protein FASTA, GFF.
  \item \textbf{Assembly stats:} custom Python: contigs from FASTA headers; GC over ATGC-only and total; GFF feature counts by column 3.
  \item \textbf{Functional gene search:} case-insensitive regex on \texttt{protein.faa} headers (arg deiminase, D/L-LDH, helveticin, bacteriocin, \texttt{tet[A-Z]}/tetracycline, \texttt{lnu}/lincosamide, \texttt{CRISPR}/\texttt{Cas}, VF-adjacent markers, mobile-element markers).
  \item \textbf{AMR / VF / plasmid screen} (the independent-tool step):
\begin{lstlisting}
for db in resfinder card ncbi argannot vfdb ecoli_vf victors plasmidfinder; do
  abricate --db $db --quiet GCA_003790365.1_ASM379036v1_genomic.fna \
    > report/evidence/abricate_${db}.tsv
done
\end{lstlisting}
    All 8 abricate DBs refreshed 2026-07-03. Cutoffs: \texttt{--minid 80 --mincov 80}.
  \item \textbf{CRISPR array search:} \texttt{minced -gffFull -spacers} on the same FASTA (MinCED v0.4).
\end{enumerate}

\section{Results vs paper}

\subsection*{C1 --- Assembly statistics}
\begin{tabular}{lrrl}
\toprule
Metric & Paper & This report & Verdict \\
\midrule
Total length (bp) & 2{,}430{,}215 & \textbf{2{,}430{,}215} & \checkmark exact \\
Contigs & 420 & \textbf{420} & \checkmark exact \\
G+C content & 39\% & \textbf{38.98\%} (ATGC-only; N=189/2.43~Mb) & \checkmark rounds to 39\% \\
\bottomrule
\end{tabular}

\subsection*{C2 --- Provenance / deposit metadata}
Illumina MiSeq \checkmark; SPAdes 3.12.0 \checkmark; piglet faeces (South Africa: Pretoria, GIT/MRS) \checkmark; PGAP annotation deposited (RAST not re-runnable free) \checkmark; deposit GCA\_003790365.1 live, GCF suppressed (``contaminated'').

\subsection*{C3 --- Named CDSs}
\begin{itemize}[leftmargin=1.2em]
  \item Arginine deiminase (EC 3.5.3.6): \textbf{present} --- ROV61345.1. \checkmark
  \item L-lactate dehydrogenase: \textbf{5 CDSs} --- ROV62718.1, ROV63569.1, ROV63627.1, ROV63895.1, ROV64206.1. \checkmark
  \item D-lactate dehydrogenase: PGAP does not use that literal name. \textbf{4 CDSs} annotated \texttt{D-2-hydroxyacid dehydrogenase} (ROV59554.1, ROV60471.1, ROV62790.1, ROV63523.1) are present --- the enzyme family that includes D-LDH; PGAP is conservative when substrate specificity is not experimentally proven. RAST (which the paper used) is more liberal. \checkmark (consistent).
  \item Bacteriocin \emph{helveticin~J}: literal ``helveticin'' not in PGAP. \textbf{1 CDS} annotated \texttt{>ROV54067.1 bacteriocin, partial} present. PGAP does not commit to a class-III subfamily; RAST does. \checkmark (consistent).
\end{itemize}

\subsection*{C4 --- Antibiotic resistance genes (``only lnuC and tetW'')}
Four independent AMR databases (ResFinder, CARD, NCBI-AMR, ARG-ANNOT), all refreshed 2026-07-03, identify \textbf{exactly the same two genes and nothing else}:

\begin{longtable}{lllrrrl}
\toprule
Database & Gene & Contig & Coords & \%Cov & \%ID & Resistance \\
\midrule
ResFinder & \texttt{tet(W)\_4} & RJWE01000089.1 & 368--2289 ($-$) & 100.00 & 99.01 & Tet/Doxy/Mino \\
ResFinder & \texttt{lnu(C)\_1} & RJWE01000125.1 & 1150--1644 ($+$) & 100.00 & 99.19 & Lincomycin \\
CARD & \texttt{tet(W)} & RJWE01000089.1 & 368--2289 ($-$) & 100.00 & 98.08 & tetracycline \\
CARD & \texttt{lnuC} & RJWE01000125.1 & 1150--1644 ($+$) & 100.00 & 99.19 & lincosamide \\
NCBI-AMR & \texttt{tet(W)} & RJWE01000089.1 & 368--2289 ($-$) & 100.00 & 99.90 & TETRACYCLINE \\
NCBI-AMR & \texttt{lnu(C)} & RJWE01000125.1 & 1150--1644 ($+$) & 100.00 & 99.19 & LINCOSAMIDE \\
ARG-ANNOT & \texttt{(Tet)tetW} & RJWE01000089.1 & 368--2289 ($-$) & 100.00 & 98.08 & --- \\
ARG-ANNOT & \texttt{(MLS)lnu(C)} & RJWE01000125.1 & 1150--1644 ($+$) & 100.00 & 99.19 & --- \\
\bottomrule
\end{longtable}

\textbf{Verdict: \checkmark exact match, quadruple-independent.} Nuance: PGAP flags this \texttt{tet(W)} locus as \texttt{pseudo=true} with \texttt{Note=frameshifted}; the abricate presence/coverage call is nevertheless correct at 100\% coverage in all four databases.

\subsection*{C5 --- Virulence factors (``no hits'')}
VFDB: \textbf{0} hits. VICTORS: \textbf{0}. ecoli\_vf: \textbf{0}. \checkmark exact match. PathogenFinder (paywalled) not rerun; with zero VF hits and only two narrow-spectrum LAB-typical AMR genes, the paper's ``zero pathogen probability'' is entirely consistent. \\
\textbf{Bonus:} PlasmidFinder detects \texttt{rep30\_1\_CDS22269(pLR581)} at 100\% cov / 100\% ID on RJWE01000056.1 --- matches the BioSample's \texttt{num\_replicons: 2}.

\subsection*{C6 --- Mobilome (prophage / IS / CRISPR)}
\begin{itemize}[leftmargin=1.2em]
  \item Prophage: 31 phage/integrase-related PGAP CDSs, including a clean structural module (phage terminase small subunit, portal, major capsid, tail tape measure, multiple tail proteins, 14 integrases). Consistent with $\geq 1$--$2$ prophages. Exact PHASTER count not rerun (paywalled). \checkmark consistent.
  \item IS elements: PGAP annotation contains dozens of transposase CDSs spanning $\geq 7$ IS families by name (IS3, IS5/IS1182, IS30, IS66, IS200/IS605, ISL3, IS21, ISLre2, IS1595). Fully consistent with paper's OASIS 7-family count. \checkmark
  \item CRISPR: MinCED found \textbf{0} arrays on this fragmented (N50 $\approx$ 28~kb, 420 contigs) assembly. Known limitation on fragmented builds. \textbf{Null-result-under-tool-choice, not a contradiction.} $\triangle$ partial.
\end{itemize}

\subsection*{C7 --- Wet-lab bacteriocin assay}
Not reproducible in silico.

\section{Verdict: REPLICATED (strong)}
\begin{itemize}[leftmargin=1.2em]
  \item 3/3 headline quantitative genome-statistic claims (C1) reproduced \textbf{exactly}.
  \item 2/2 provenance claims (C2) reproduced exactly.
  \item 3/3 named-CDS classes of C3 reproduced (D-LDH via PGAP family name; bacteriocin by generic name).
  \item \textbf{C4 (the paper's key safety claim) reproduced 4$\times$ independently} with identical hits, coordinates, and near-identical \%ID --- the strongest single result.
  \item \textbf{C5 (no virulence factors) reproduced 3$\times$ independently.}
  \item C6 partially reproduced.
  \item C7 not reproducible in silico.
\end{itemize}

\textbf{No finding here contradicts the paper.}

\section{Genuine Critique}
This section is deliberately adversarial, in the spirit of what a critical peer reviewer or replication-integrity auditor would flag. None of these observations negate the REPLICATED verdict, but they qualify how the paper's headline framing should be read.

\begin{enumerate}[leftmargin=1.6em]

  \item \textbf{The RefSeq mirror is suppressed for ``contamination'' --- the paper does not mention this.} \\
  NCBI's automated QC flagged the assembly \texttt{warnings: ["contaminated"]}, and RefSeq staff removed GCF\_003790365.1 on those grounds (``This record was removed by RefSeq staff... Reason: contaminated''). The paper does not discuss this at all. This does not change the AMR/VF calls (both target well-characterised loci with high sequence identity that survive most contamination scenarios), but it does mean the \emph{fine-grained gene inventory} --- specifically claims about the absence of unexpected genes --- should be treated with more caution than the paper acknowledges. Contamination inflates apparent gene content; absence claims are correspondingly weaker.

  \item \textbf{Assembly is highly fragmented; some paper-level claims are structurally unverifiable on this deposit.} \\
  N50 = 28{,}048~bp, L50 = 24, 420 contigs. That is what the paper reports and what I confirm --- but it means CRISPR array detection (MinCED) legitimately failed here, and the paper's ``5 CRISPR CDSs each with Cas'' claim depends on the specific tool (CRISPRFinder) and the specific RAST-annotated build. On a differently annotated deposit of the \emph{same} raw reads, CRISPR content might differ. Independent CRISPR detection on the deposit is therefore \textbf{null under my tool choice, not concordant}. Fair, but a soft spot.

  \item \textbf{Annotator identity matters more than the paper implies.} \\
  Two ``present in the paper, not literally present in my PGAP re-check'' cases --- helveticin~J and D-lactate dehydrogenase (by name) --- are both attributable to RAST vs PGAP naming conservatism. This is a scientifically legitimate difference, but the paper does not distinguish which annotator produced which claim. If a downstream user re-annotates with a different tool (Prokka, Bakta), they will not reproduce the literal names either; they will reproduce the underlying biology.

  \item \textbf{The \texttt{tet(W)} locus is a PGAP-annotated pseudogene (\texttt{frameshifted}); ``presence'' $\ne$ functional resistance.} \\
  The paper treats \texttt{tet(W)} as a resistance gene (presence-based screen). PGAP flags this specific locus as \texttt{pseudo=true, Note=frameshifted}. Every abricate database still calls it at 100\% coverage (because sequence homology is intact), so the paper's methodology --- a homology-based screen --- is not wrong. But the biological implication (does this strain confer tetracycline resistance?) is weaker than the paper implies. A phenotypic MIC panel would settle it; the paper did not do that for the AMR calls (only for the bacteriocin activity).

  \item \textbf{PathogenFinder was not independently re-run; the ``zero pathogen probability'' remains single-source.} \\
  My replication confirms VFDB / VICTORS / ecoli\_vf = 0 hits (3$\times$ agreement). But PathogenFinder is a CGE hosted service with its own model, and I did not rerun it. The paper's zero-probability call is thus \emph{consistent} with my independent VF screen but not \emph{independently reproduced by the same tool}.

  \item \textbf{C6 is the weakest triangulated claim.} \\
  Two of its three sub-claims (prophage count, CRISPR count) rest on paywalled or fragmented-assembly-sensitive tools. Only the IS-family count is robustly triangulated here.

  \item \textbf{The wet-lab claim (C7) is unreachable in silico by construction.} \\
  Zone-of-inhibition against STEC O177 with a partially purified bacteriocin preparation is a laboratory result. It should not be counted as ``replicated'' or ``not replicated'' by any in silico exercise; I have flagged it as \emph{not reproducible in silico}, which is the honest label.

  \item \textbf{Generalisability: one deposit, one annotator, one day of tool state.} \\
  All abricate DBs were refreshed 2026-07-03 (same day). This is good for reproducibility snapshotting but means a rerun in 6 months against updated databases may find new hits (particularly if ResFinder / CARD add new tet or lnu variants, or new VF entries). The AMR/VF landscape is not static.

\end{enumerate}

\noindent \textbf{Net.} The paper's safety framing (``only two narrow-spectrum AMR genes; no virulence factors'') is well supported by independent evidence at the level the paper actually screens. The paper's higher-level claim (``high-quality reference for a probiotic candidate'') is somewhat weakened by the RefSeq-suppression / fragmentation issues that the paper does not discuss. The scientifically important claims replicate; the presentational claims deserve a caveat.

\section{Reproducibility footprint}
All raw evidence in \texttt{report/evidence/}:
\begin{itemize}[leftmargin=1.2em, itemsep=0pt]
  \item \texttt{ncbi\_datasets\_report.json} --- strain metadata (2 records for PNW1)
  \item \texttt{assembly\_stats.json} --- contigs, total\_bp, GC\%, N-bases, GFF feature counts
  \item \texttt{gene\_search.json} --- regex hits in \texttt{protein.faa}
  \item \texttt{abricate\_\{resfinder,card,ncbi,argannot\}.tsv} --- AMR (C4)
  \item \texttt{abricate\_\{vfdb,victors,ecoli\_vf\}.tsv} --- VF (C5)
  \item \texttt{abricate\_plasmidfinder.tsv} --- bonus plasmid rep-protein
  \item \texttt{minced\_crispr.gff} --- CRISPR (empty output)
\end{itemize}

Total wall time on CherryRd (no GPU, no paid API): $\sim$3 minutes. Zero paid endpoints. All tools FOSS; all data open NCBI GenBank.

\end{document}
